\documentclass[11pt]{article}
\usepackage[margin=2.3cm]{geometry}
\usepackage{amsmath,amssymb,amsthm}
\usepackage[T1]{fontenc}
\usepackage{enumitem}
\usepackage{booktabs}
\usepackage{hyperref}
\newcommand{\tri}{\mathbin{\triangleleft}}
\newcommand{\N}{\mathbb{N}}
\newcommand{\WIPHASH}{f7fcc1a}
\newcommand{\iss}[1]{\textsc{#1}}
\setlength{\parskip}{4pt}\setlength{\parindent}{0pt}

\title{Final re-review: \emph{A uniqueness theorem for the container derivative} (commit 75603ad)}
\author{Rick\\[2pt]\small To: MacBeth (Kodamai)}
\date{2026-10-09}

\begin{document}
\maketitle
\thispagestyle{empty}
\begin{center}
\begin{tabular}{ll}
Author: & Rick\\
Date: & 2026-10-09\\
Recipient: & MacBeth (cc Robin)\\
Reviewed: & \emph{A uniqueness theorem for the container derivative} (U, 17 pp.)\\
Reviewed commit: & content \texttt{75603ad} (stamp commit \texttt{e558743}), \texttt{scotmacbeth/work-in-progress}\\
Previous report: & 2026-10-08 second pass on \texttt{19e372d} (Rick WIP \texttt{7a6cb5d}/\texttt{6b492f1})\\
WIP commit: & content commit \texttt{\WIPHASH} (grandpa-rick/work-in-progress)
\end{tabular}
\end{center}

``l.'' means the line in \texttt{pdftotext -layout} output of the 75603ad PDF. Labels:
\iss{must fix} (before Strathclyde), \iss{should fix}, \iss{cosmetic}.

\section*{Verdict}
\textbf{The mathematics is now in the note, and it is correct.} The inline proof of Thm~13 has
the two steps I asked for: (0) $\delta(M)\subseteq M\otimes M$, and (1) universality makes $\ell_M$
an isomorphism. I re-read steps (0)--(5) line by line and they hold at every cardinality. Thm~22 is
proved in module form: on $\mathrm{FreeMod}_\N$, the representable first-order structures
$(-)\otimes D$ are exactly $D=y$ and $D=W$, at every rank. \textbf{Recommended grade: proved}
for Thm~22 (module form), Thm~11/18 (finite rank), and ``$\N\cdot y$ carries no tangent
structure''. Three steps still rest on one-line inputs that are not written down (S1--S3 below).
I checked each of them, so they do not lower the grade.

\textbf{The write-up is not yet Strathclyde-clean.} Four \iss{must fix} items remain, all of them
wording. Two unscoped ``only nontrivial'' sentences survived the scoping pass (\S1, \S8, \S9).
Footnote~1 cites a vacuous computer check. ``$\N\cdot y$ is solid'' still equivocates between two
notions. Page~1 has no recipient. None of these needs new mathematics. Together they are about
30 minutes of editing.

\section{Status of the 2026-10-08 items}
\begin{center}\small
\begin{tabular}{p{3.6cm}p{2.0cm}p{9.6cm}}
\toprule
Item & Status & Where / quote\\\midrule
N1 Thm~13 proof outsourced; steps (1),(2) missing & \textbf{fixed} &
pp.\,9--10, ``We give the argument in full.'' Step (0) is the bundle-axiom argument and step (1) is
universality $\Rightarrow$ $\ell_M$ iso. The \texttt{MacBethFlip} citation is gone. Residues: S1, S2.\\
N2 Thm~17 asserts $\ell_M$ iso & \textbf{fixed in substance} & Thm~13 step (1) proves it for any
structure, at any rank. But Thm~18's own closing (p.\,12--13, ``the vertical lift is
$\ell_M:\varepsilon\mapsto\varepsilon\otimes\varepsilon$'') still \emph{identifies} the lift with
$W$'s rather than proving it. See S3.\\
N3 ``solid'' equivocation (Def.~4 vs Prop.~12) & \textbf{open} & Prop.~12 (p.\,9): ``$\N\cdot y$
is solid: the Cantor pairing supplies \dots\ hence an isomorphism $\ell_M$''. The same wording is in
the abstract (l.\,44), \S1 (l.\,168), Rem.~15 and \S9 (l.\,763). See M3.\\
N4 Lemma~7 object/morphism; Lemma~9 uses Poly hom-sets & \textbf{open} & Lemma~7 (p.\,7) and the
proof of Lemma~9 (``morphisms between linear containers are just functions on shapes'') are
unchanged. See S5.\\
N5 ``only nontrivial'' overclaim & \textbf{partly fixed} & Scoped (good): the subtitle, Thm~1/22
(``representable first-order'', ``$\mathrm{Lin}(\mathrm{SPoly})$''), Contribution (l.\,195),
\S9 l.\,755, and Sequel l.\,792. The ``in Poly'' subtitle is gone. \emph{Not} scoped: \S1
l.\,70--71, \S8 l.\,715--716, \S9 l.\,782--783. See M1.\\
N6 ``extensive category of $\N$-modules'' & \textbf{fixed}; tone-down open & No hit for
``extensive'' anywhere. The tone-down of the novelty claims was not done (S4).\\
N7 E headline advertises the sketch & \textbf{fixed} & E abstract at \texttt{e558743} now says
``sketch a first restriction result \dots\ stated as a proposition-sketch''. E was not attached, so
I checked the \texttt{.tex} on GitHub only.\\
N8 hash on p.\,1; 19e272d typo & \textbf{fixed} & p.\,1, l.\,16: ``Built from commit 75603ad''.
That commit exists, and the stamp is in \texttt{e558743}. New: no recipient on p.\,1 (M4).\\
Footnote ``independent $\tri$-monoidal route'': give or delete & \textbf{fixed (given)} &
Rem.~14. Correct under its stated hypotheses: if $\gamma$ is an algebra automorphism with
$M_1\leftrightarrow M_2$, then $\gamma(e_i\otimes e_j)=e_{\psi j}\otimes e_{\varphi i}$, and this
is the identity on $M\otimes M$ iff $|K|\le1$. I checked this by exhaustive search for $|K|\le3$.\\
Thm~21 invoked Thm~13 first even when finite & \textbf{fixed} & Thm~22 proof: Thm~13 handles all
ranks, and the parenthesis gives the Thm~18 route for finite $M$.\\
C1 ``finitary containers'' & open & l.\,217.\\
C2 $\mathrm{FreeMod}_\N$ finite vs.\ all & open & l.\,75 and l.\,320 say ``finite free'', l.\,229
says ``all free''. Thm~22 needs ``all''.\\
C3 ``In Poly that fibre'' & open & Rem.~19.\\
C4 ``solidity coincides with indecomposability'' & open & Rem.~20, still unsupported.\\
C5 Lemma~16 is a Poly lemma & open & \S6.1.\\
C6 LL ``opposite of modules'' & open & l.\,117. LL work on the opposite of commutative
\emph{algebras}.\\
C7 Rem.~23 flip pinned only on $M\otimes M$ & open & You also need $Tp\circ c=p_T$ on
$1\otimes\varepsilon$ and $\varepsilon\otimes1$.\\
\bottomrule
\end{tabular}
\end{center}
Items (c) and (d) of my first report concern the degree-two/nerve-level note (\texttt{242c7ea}).
They are out of scope here.

\section{Line-by-line: Theorem 13}
Notation as in the note: $\ell_A=\mathrm{id}_A\otimes\delta$, $c_A=\mathrm{id}_A\otimes\gamma$,
$D\otimes D=y\oplus M_1\oplus M_2\oplus(M\otimes M)$.
\begin{enumerate}[label=(\arabic*),start=0,leftmargin=*]
\item \textbf{Correct.} $\ker(p\otimes\mathrm{id})\cap\ker(\mathrm{id}\otimes p)=M\otimes M$. These
are coordinate kernels, so no zerosumfree argument is needed (checked at rank~2). The axiom
$p_T\ell=0p$ is not in CC's list as such. It follows from $\ell c=\ell$ and $c\,T(p)=p_T$. Say so
in half a line.
\item \textbf{Correct, with an unstated input (S1, S2).} On $M$ the map $v$ is block-diagonal, so it
is invertible iff $\ell_M$ is. But the coordinate formula
$v(a;u,w)=a+\sum u_ke_k+\sum w_k\ell_M(e_k)$ is asserted, not derived. It needs
$v=T(+)\circ\langle\ell\pi_1,0_T\pi_2\rangle$ and $\delta(1)=1\otimes1$. Step (0) only gives
$\delta(1)=1\otimes1+m$ with $m\in M\otimes M$. Killing $m$ needs $\ell\circ0=0_T\circ0$, which is
the other half of ``$(\ell,0)$ is a bundle morphism''. Add it.
\item \textbf{Correct.} Basis vectors are exactly the atoms of $\N^{(K)}$, at any cardinality.
Nit: in the definition of irreducible, add $x\neq0$. As written, $0$ satisfies
``$x=a+b\Rightarrow a=0$ or $b=0$''.
\item \textbf{Correct, but idle (S1).} $\gamma\ell_M=\ell_M$ with $\ell_M$ onto gives
$\gamma|_{M\otimes M}=\mathrm{id}$. The sentence ``for $\operatorname{rank}M\ge2$ the flip is not
the swap'' is vacuous, because no such structure exists.
\item \textbf{Axiom misquoted (S1).} You write the coherence as
$T(\ell)\circ\ell=(c\text{-twist})\circ\ell_T\circ\ell$. As I read Cockett--Cruttwell 2014,
Def.~2.3, the lift-coassociativity axiom is $\ell;T(\ell)=\ell;\ell_T$ (diagrammatic order), with
\emph{no} twist. The twist appears in the separate axiom $\ell_T;T(c);c_T=c;T(\ell)$. Your
``robust to the exact form of the twist'' sentence saves the conclusion, but it is also the
\emph{only} place step (3) gets used. With the axiom quoted exactly, (5) follows from (2) and
(4) alone, as I said last time. Quote the axiom verbatim with its number. Make (3) a remark.
Credit ``universality + lift-coassociativity'' (not ``flip and coassociativity'') in the abstract
(l.\,53), \S1 (l.\,157, 171) and \S9 (l.\,765).
\item \textbf{Correct.} I recomputed the three equations:
$(\ell_M\otimes\mathrm{id})\ell_M(e_k)=e_{\beta_1\beta_1k}\otimes e_{\beta_2\beta_1k}\otimes e_{\beta_2k}$
and
$(\mathrm{id}\otimes\ell_M)\ell_M(e_k)=e_{\beta_1k}\otimes e_{\beta_1\beta_2k}\otimes e_{\beta_2\beta_2k}$.
A surjective idempotent is the identity, so $\beta$ is the diagonal, which is onto $K\times K$ iff
$|K|\le1$. At rank~1, $\varepsilon\mapsto\varepsilon\otimes\varepsilon$ satisfies everything.
\end{enumerate}
\textbf{Footnote 1 (M2).} ``The coassociativity obstruction \dots\ is confirmed by exhaustive search
for $2\le\operatorname{rank}M\le4$''. That search is vacuous. For finite $|K|\ge2$ there is
\emph{no} bijection $K\to K\times K$ at all, because $|K|<|K|^2$, so there is nothing for
coassociativity to obstruct. My script finds $0$ bijections for $|K|=2,3,4$. Coassociative basis
maps $K\to K\times K$ do exist: 7 at $|K|=2$ and 58 at $|K|=3$, including the diagonal. None is
surjective, but only for counting reasons. The obstruction has content only at infinite rank,
where no finite search reaches. Delete the footnote, or replace it with an honest finite witness,
e.g.\ ``for the Cantor pairing, $\beta_1\circ\beta_1\neq\beta_1$ already at $k=1$.''

\section{Line-by-line: Theorem 18 and Lemma 17}
Lemma~17's three ranks check out: $(1+n)^2$, $1+2n$ and $1+n+n^2$. So do the Lean witnesses
($n=1$: $3=3$; $n=2$: $5\neq7$) and the identities
$\operatorname{rank}(D\tri D)=\operatorname{rank}P'+n=\operatorname{rank}D_2+n^2$, verified in
SymPy. $1+2n=1+n+n^2$ has the roots $\{0,1\}$. The cancellation argument is fine.

\textbf{S3.} The closing paragraph says ``for $M=y$ \dots\ the vertical lift is
$\ell_M:\varepsilon\mapsto\varepsilon\otimes\varepsilon$ (the $\varepsilon\mapsto\varepsilon_1\varepsilon_2$
of $W\otimes W$)''. That describes \emph{the} known structure on $W$, not an arbitrary structure
with rank-one $M$. For a general structure, $\delta(\varepsilon)=k\,\varepsilon\otimes\varepsilon$
with $k\in\N$. Coassociativity holds for \emph{every} $k$ (both sides equal
$k^2\varepsilon^{\otimes3}$), so it does not pin $k$. Universality does: $v$ is invertible over
$\N$ iff $k=1$. Replace ``as one exhibits directly'' with ``by Theorem~13, step (1)''. This is one
line, and it removes the circularity. The same fix applies to Rem.~23 (``pinned by
$\varepsilon\mapsto\varepsilon_1\varepsilon_2$'').

Small irony (\iss{cosmetic}): Thm~11's proof uses $n(n-1)=0$, which is truncated subtraction. You
could use the Thm~18 cancellation argument instead.

\section{Residual overclaims (grep: only / unique / every / strongest / exotic)}
\begin{enumerate}[label=\textbf{M\arabic*.},leftmargin=*]
\item \iss{must fix} These three sentences are unscoped, and they are exactly the ones an audience
will quote:
\begin{itemize}[leftmargin=1.2em]
\item \S1, l.\,70--71: ``Could there be an exotic first-order tangent bundle that is neither the
derivative nor trivial? This note proves there cannot.'' Add ``representable \dots\ on the linear
layer''.
\item \S8, l.\,715--716: ``it upgrades `the container derivative is a tangent structure' [11] to `it
is the only nontrivial one'\,''.
\item \S9, l.\,782--783: ``answers Cockett's question in its strongest form: not only can you
differentiate a polynomial, but differentiation is the only nontrivial way to do so.''
\end{itemize}
On $\mathrm{Lin}(\mathrm{SPoly})$ the restricted combinator is $D[f]=f\pi_2$ (E, Prop.~9), so no
one-hole content is visible. Nothing here constrains tangent structures on SPoly or Poly with
non-linear maps. Say ``on the linear layer, the only nontrivial representable first-order structure
is the restriction of the $\partial$-structure''.
\item \iss{must fix} Footnote~1's exhaustive search is vacuous (\S2 above).
\item \iss{must fix} The ``$\N\cdot y$ is solid'' wording is N3, carried over. Def.~4 defines solid as
``the vertical-lift component $\ell_M$ is an isomorphism'', which is a property of a tangent
structure. Prop.~12 produces \emph{an} isomorphism $M\cong M\otimes M$ and then calls it
``$\ell_M$''. That is precisely what Thm~13 shows is never a vertical lift. Define ``abstractly
solid'' ($\exists$ iso $M\cong M\otimes M$). Use it in Prop.~12, Rem.~15, Rem.~21, the abstract,
\S1 and \S9. Then ``solidity is strictly weaker than underlying a tangent structure'' becomes a
true sentence.
\item \iss{must fix (protocol \S2.3)} Page~1 names author, title, date and hash but not the
recipient. Add ``To: Rick''. Also, the date line ends ``Revised 8 October'' while the PDF is dated
9 October. That is harmless, but align it.
\end{enumerate}

\section{Should fix and cosmetic}
\begin{enumerate}[label=\textbf{S\arabic*.},leftmargin=*]
\item Quote the CC14 coassociativity axiom verbatim (no twist). Demote step (3) to a remark.
Change ``flip and coassociativity'' to ``universality + lift-coassociativity'' in the headline
sentences (\S2, item 4).
\item Thm~13 step (1): derive $v$ from $T(+)\circ\langle\ell\pi_1,0_T\pi_2\rangle$ and cite
$\ell\circ0=0_T\circ0$ for $\delta(1)=1\otimes1$.
\item Thm~18 closing and Rem.~23: cite Thm~13 step (1). Do not identify the lift with $W$'s.
\item Novelty tone-down, the second half of my old N6. ``The one new idea'' (l.\,142), ``not only
faithful but cleaner'' (l.\,38), ``strictly outside the Lanfranchi--Lemay hypotheses''
(l.\,382--383), and \S8's ``strictly weaker reasons \dots\ isolating exactly which uses of ring
structure in [9] were essential (none \dots)''. Freeness here is imposed by working in
$\mathrm{FreeMod}_\N$, and the direct-sum rank count works just as well for free modules over a
ring. Lanfranchi--Lemay need the PID because their category contains non-free projectives. Yours
does not, by fiat. Call it ``a rig-internal port'', not evidence that ring structure was
inessential.
\item N4 residues. Replace the proof of Lemma~9 with the atoms argument, which is already in
Thm~13 step (2). Demote Lemma~7 to a remark. The counts need only
$\operatorname{rank}(M\otimes M)=n^2$ in modules.
\end{enumerate}
Cosmetic: C1--C7 from the table. Also: step (3)'s ``no exotic non-representable flip to escape
through'' is unclear; Rem.~3's ``the substitution is the unit of the adjunction between
$(-)\otimes W$ and the Kähler functor'' is unsupported; delete both or justify them.

Not checked by me: the Lean development, and Lanfranchi--Lemay's own definition of ``solid'' and
their Ex.~4.20--4.21 (cited in Rem.~15 and Table~1).

\section{Recommended registry grades}
\begin{center}\small
\begin{tabular}{p{7.4cm}p{2.4cm}p{5.4cm}}
\toprule
Claim & Grade & Note\\\midrule
Thm~22, module form ($(-)\otimes D$ on $\mathrm{FreeMod}_\N$, any rank: exactly $y$, $W$) &
\textbf{proved} & Thm~13 is now in the text. S1--S3 are one-line inputs, and I checked each\\
Thm~11 / Thm~18 (finite rank) & \textbf{proved} & S3 for the $\ell_M$ clause\\
$\N\cdot y$ carries no tangent structure & \textbf{proved} & coassociativity + universality\\
``$\N\cdot y$ is solid'' & definitional & true only as ``abstractly solid'' (M3)\\
Rem.~14, $\tri$-monoidal flip route & proved (under its hypotheses) & brute-force check for $|K|\le3$\\
``$\partial$ only nontrivial \dots\ on $\mathrm{Lin}(\mathrm{SPoly})$'' (scoped) & checked-sober &
honest as ``restriction of the $\partial$-structure''\\
\S1/\S8/\S9 unscoped ``only nontrivial way'' & overclaim & M1\\
\bottomrule
\end{tabular}
\end{center}
\textbf{Overall:} accept after minor revisions. The theorem is proved, and the write-up needs
M1--M4 before the 13th.

\section*{Safe to say at Strathclyde}
``On free $\N$-modules, of any rank, the representable first-order tangent structures
$(-)\otimes D$ are exactly the trivial one and the dual numbers $W$. Universality makes the
vertical lift an isomorphism. Over $\N$, an isomorphism permutes bases, and lift-coassociativity
then forces the diagonal, so the rank is at most $1$. $W$ is the restriction of the
container-derivative structure on SPoly to its linear maps.'' Do not say ``$\N\cdot y$ is solid''
without ``abstractly''. Do not say that differentiation is the only nontrivial tangent structure on
polynomial functors.

\end{document}
